\documentclass[main]{subfiles}

\begin{document}

\chapter{Foldy--Wouthuysen Derivations and the SME--Dobrescu--Mocioiu Mapping}

\section{Overview}

Chapter~2 introduced the Standard Model Extension (SME) fermion Lagrangian in
full generality, the complete Dobrescu--Mocioiu (DM) potential catalogue
$V_1$--$V_{16}$, and the general Foldy--Wouthuysen (FW) transformation
machinery. This chapter applies that machinery to the three SME coefficients
most relevant to spin-dependent exotic interactions -- $b_\mu$, $H_{\mu\nu}$,
and $d_{\mu\nu}$ -- deriving the non-relativistic Hamiltonian each one
generates and matching it onto the corresponding DM potential(s) from
Chapter~2's catalogue. Section~\ref{sec:bmu} treats $b_\mu$, the CPT-odd
coefficient responsible for a genuine matter--antimatter sign flip;
Section~\ref{sec:Hmunu} treats $H_{\mu\nu}$, a CPT-even tensor coupling;
Section~\ref{sec:dmunu} treats $d_{\mu\nu}$, the CPT-even axial-tensor
partner of $H_{\mu\nu}$, whose $d_{i0}$ derivation is flagged as unverified
pending further work. Section~\ref{sec:master} then synthesises all three
derivations into a master SME$\to$DM matching table and the interpretive
rule connecting predicted CPT parity to the asymmetry parameter $A_\alpha$
actually computed by SPINDEP.

A recurring theme in this chapter is worth stating up front: $A_\alpha$, as
computed from experimental upper bounds, cannot by itself distinguish a
genuine CPT-violating signal from an ordinary difference in experimental
sensitivity between the matter- and antimatter-sector measurements being
compared. This is derived carefully in Section~\ref{sec:bmu-asymmetry} and
then applied uniformly to every SME coefficient in Section~\ref{sec:master}.

\section{Foldy--Wouthuysen Derivation of the \texorpdfstring{$b_\mu$}{b\_mu} Term}
\label{sec:bmu}

\subsection{Overview and Physical Motivation}

The SME coefficient $b_\mu$ (sometimes written $b_\mu$ or $b_\mu^f$ for
fermion $f$) parametrises CPT-odd, Lorentz-violating interactions of the
form $-b_\mu\bar\psi\gamma_5\gamma^\mu\psi$ in the SME Lagrangian. This term
is CPT-odd because, under CPT, $\psi\to i\gamma^0\gamma^2\psi^*$ and the
combination $\bar\psi\gamma_5\gamma^\mu\psi$ changes sign. It is the leading
CPT-violating operator relevant to spin-dependent exotic interactions because
it couples to the fermion spin directly through the $\gamma_5$ structure.

The goal of this derivation is to take the fully relativistic SME Lagrangian
term and obtain the non-relativistic (NR) Hamiltonian via the FW
transformation, then match the NR Hamiltonian to the DM catalogue to identify
which potential(s) $V_n$ it generates in a two-body interaction.

\paragraph{Key result.} The $b_i$ (spatial) component generates $V_2$ at
leading order $m^0$: $H_{NR}=-\mathbf b\cdot\boldsymbol\sigma$ for matter,
$+\mathbf b\cdot\boldsymbol\sigma$ for antimatter (Section~\ref{sec:bmu-cconj}).
The $b_0$ (temporal) component generates $V_7,V_8$ at subleading order
$m^{-1}$: $H_{NR}=+b_0(\boldsymbol\sigma\cdot\mathbf p)/m$. This sign flip is
the qualitative CPT-odd signature of $b_\mu$; it does \emph{not} by itself
predict $|A_\alpha|\to 1$ in SPINDEP's output, for reasons developed in
Section~\ref{sec:bmu-asymmetry}.

\subsection{The SME Lagrangian Term}

\subsubsection{Covariant Form}

The CPT-odd Lorentz-violating modification to the free Dirac Lagrangian in
the minimal SME is
\begin{equation}
  \mathcal{L}_b = -b_\mu\,\bar\psi\,\gamma_5\gamma^\mu\,\psi. \label{eq:Lb}
\end{equation}
The corresponding modified Dirac equation is
\begin{equation}
  \left(i\gamma^\mu\partial_\mu - m - b_\mu\gamma_5\gamma^\mu\right)\psi = 0.
\end{equation}

\subsubsection{CPT Properties}

Under the CPT transformation, the spinor transforms as
$\psi\to i\gamma^0\gamma^2\psi^*_{\text{CPT}}$ (up to a phase), and the
relevant bilinear transforms as
\begin{equation}
  \mathrm{CPT}\!\left[\bar\psi\gamma_5\gamma^\mu\psi\right]
  = -\bar\psi\gamma_5\gamma^\mu\psi.
\end{equation}
So $\mathcal L_b\to+b_\mu\bar\psi\gamma_5\gamma^\mu\psi$ under CPT, which
differs from the original $-b_\mu\bar\psi\gamma_5\gamma^\mu\psi$ by a sign,
confirming that $b_\mu$ is CPT-odd. Under charge conjugation $C$ alone
(mapping particles to antiparticles),
\begin{equation}
  C\!\left[\bar\psi\gamma_5\gamma^\mu\psi\right]
  = +\bar\psi\gamma_5\gamma^\mu\psi
  \qquad (C\text{ is even for the axial current}),
\end{equation}
but the \emph{coupling} changes sign because the antiparticle's effective
$b_\mu$ coefficient is $-b_\mu$ relative to the particle's. This is the
origin of the matter--antimatter sign flip derived explicitly in
Section~\ref{sec:bmu-cconj}.

\subsection{Application of the Foldy--Wouthuysen Transformation}

This section applies the general FW expansion derived in Chapter~2 --
reproduced here for convenience as
\begin{equation}
  H' \approx \beta m + \varepsilon_{\text{even}}
  + \frac{\beta\varepsilon_{\text{odd}}^2}{2m}
  - \frac{[\varepsilon_{\text{odd}},\varepsilon_{\text{even}}]}{4m}
  + O(m^{-2}), \label{eq:FWexpansion}
\end{equation}
with $H=\beta m+\varepsilon_{\text{even}}+\varepsilon_{\text{odd}}$,
$\beta=\gamma^0$, and $\varepsilon_{\text{even}},\varepsilon_{\text{odd}}$
commuting and anticommuting with $\beta$ respectively -- to the specific
case of the $b_\mu$ coefficient.

\subsubsection{Application to the \texorpdfstring{$b_\mu$}{b\_mu} Term}
\label{sec:bmu-application}

The $b_\mu$ Hamiltonian, obtained from Eq.~\eqref{eq:Lb}, is
\begin{equation}
  H_b = b_\mu\gamma^0\gamma_5\gamma^\mu
      = b_0\gamma^0\gamma_5\gamma^0 + b_i\gamma^0\gamma_5\gamma^i.
\end{equation}

\paragraph{Spatial components $b_i$.} In the Dirac representation, where
$\gamma^0=\mathrm{diag}(1,1,-1,-1)$ and $\gamma_5$ is off-diagonal, the term
$\gamma^0\gamma_5\gamma^i$ is an even operator (block-diagonal). Using
$\{\gamma_5,\gamma^i\}=0$, the operator identity is
\begin{equation}
  \gamma^0\gamma_5\gamma^i = -\Sigma^i = -\,\mathrm{diag}(\sigma^i,\sigma^i),
\end{equation}
where $\Sigma^i$ is the $4\times4$ spin matrix. In the upper $2\times2$
block, this gives directly
\begin{equation}
  \boxed{H_{NR}(b_i) = -\mathbf b\cdot\boldsymbol\sigma \quad (\text{matter})}
  \label{eq:HNR-bi-matter}
\end{equation}
already in Pauli non-relativistic form, requiring no further FW iteration at
leading order in $1/m$: the $b_i$ term is exact at order $m^0$ in the NR
expansion.

This result depends on the ordering of $\gamma_5$ and $\gamma^\mu$ in the
Lagrangian bilinear: $\bar\psi\gamma_5\gamma^\mu\psi$ and
$\bar\psi\gamma^\mu\gamma_5\psi$ are exact negatives of each other, since
$\{\gamma_5,\gamma^\mu\}=0$. This derivation follows
$\mathcal L_b=-b_\mu\bar\psi\gamma_5\gamma^\mu\psi$
(Eq.~\eqref{eq:Lb}), matching the standard SME mass-operator form
$M=m+a_\mu\gamma^\mu+b_\mu\gamma_5\gamma^\mu+\tfrac12H_{\mu\nu}\sigma^{\mu\nu}$
used throughout Kostelecký's papers, including Kostelecký and Lane (1999),
which is cited throughout this chapter.

\paragraph{Temporal component $b_0$.} The temporal component produces an odd
operator,
\begin{equation}
  \gamma^0\gamma_5\gamma^0 = \gamma_5 \quad (\text{off-diagonal}),
\end{equation}
which does not contribute at order $m^0$. At order $m^{-1}$, using the FW
expansion of Eq.~\eqref{eq:FWexpansion}, the odd operator generates
\begin{equation}
  H_{NR}(b_0) = \frac{\beta(\gamma_5 b_0)^2}{2m} + \ldots
  \;\longrightarrow\; \frac{b_0(\boldsymbol\sigma\cdot\mathbf p)}{m}
  + O(m^{-2}),
\end{equation}
which is velocity-dependent and matches DM potentials $V_7$ and $V_8$.

\subsection{Charge Conjugation and the Antimatter Sign Flip}
\label{sec:bmu-cconj}

\subsubsection{C-transformation of the Spinor}

Under charge conjugation, the Dirac spinor transforms as
\begin{equation}
  \psi \to \psi^c = C\gamma^0\psi^*, \qquad C = i\gamma^2\gamma^0
  \quad (\text{Dirac convention}).
\end{equation}
The charge-conjugated spinor describes the antiparticle with the same
momentum but opposite charge. Crucially, the SME coefficient $b_\mu$ is a
fixed background field that does \emph{not} transform under $C$ -- it is an
external source -- so the Lagrangian density for the antiparticle becomes
\begin{equation}
  \mathcal L_b(\psi^c) = +b_\mu\,\bar\psi^c\gamma_5\gamma^\mu\psi^c
  = +b_\mu\left(\bar\psi\gamma_5\gamma^\mu\psi\right)^*.
\end{equation}
The sign flip from $-b_\mu$ (particle) to $+b_\mu$ (antiparticle) follows
from the anticommutativity of the $C$ matrix with $\gamma_5$:
\begin{equation}
  C\gamma_5\gamma^\mu C^{-1} = -(\gamma_5\gamma^\mu)^*
  \;\Longrightarrow\;
  \bar\psi^c\gamma_5\gamma^\mu\psi^c = -\left(\bar\psi\gamma_5\gamma^\mu\psi\right).
\end{equation}
Combined with the $-b_\mu$ in the Lagrangian, the antiparticle coupling has
overall $+b_\mu$, so
\begin{equation}
  H_{NR}^{\bar f}(b_i) = +\mathbf b\cdot\boldsymbol\sigma. \label{eq:HNR-bi-anti}
\end{equation}
Equations~\eqref{eq:HNR-bi-matter} and \eqref{eq:HNR-bi-anti} together give
the genuine CPT-odd sign flip for $b_i$: $-\mathbf b\cdot\boldsymbol\sigma$
for matter, $+\mathbf b\cdot\boldsymbol\sigma$ for antimatter.

\subsubsection{Implication for the Asymmetry Parameter}
\label{sec:bmu-asymmetry}

The two-body interaction potential for a matter--antimatter pair is
proportional to the product of the single-particle NR Hamiltonians. For the
$b_i$ term,
\begin{equation}
  V_{\text{pair}} \propto H_{NR}^{f}\times H_{NR}^{\bar f}
  \propto (-\mathbf b\cdot\boldsymbol\sigma_1)(+\mathbf b\cdot\boldsymbol\sigma_2),
\end{equation}
a potential that changes sign relative to the matter--matter case, where
both factors carry $-\mathbf b\cdot\boldsymbol\sigma$. The coupling constant
$g$ measured from matter--matter and matter--antimatter experiments
therefore differ by a sign, yielding $g_m=-g_{\bar a}$ \emph{for a
hypothetical, exactly known, signed coupling}. Substituting into the SPINDEP
asymmetry parameter,
\begin{equation}
  A_\alpha = \frac{g_m-g_{\bar a}}{g_m+g_{\bar a}}
  = \frac{g-(-g)}{g+(-g)} = \frac{2g}{0} \;\longrightarrow\; \infty.
  \label{eq:Aalpha-divergence}
\end{equation}
This is a genuine divergence -- the denominator vanishes identically -- not a
value that saturates at $1$. An exact CPT-odd sign flip of a \emph{signed}
coupling makes Eq.~\eqref{eq:Aalpha-divergence} formally undefined, not
bounded; this was verified directly with a symbolic limit calculation
(\texttt{sympy.limit}) in the accompanying derivation notebooks.

The resolution is that Eq.~\eqref{eq:Aalpha-divergence} does not describe
what SPINDEP actually computes. SPINDEP's $g_m/g_{\bar a}$ inputs are
\emph{always-positive experimental upper bounds} from independent
experiments, not a signed measurement of one real underlying coupling. For
any two independent positive bounds $g_1,g_2$, the ratio
$(g_1-g_2)/(g_1+g_2)$ approaches $\pm1$ whenever one bound is far tighter
than the other, \emph{regardless of whether the underlying physics is
CPT-odd or CPT-even}. This is a generic property of comparing two positive
numbers of very different size -- a sensitivity-gap effect -- not a
signature of the $b_\mu$ mechanism specifically: substituting
$g_{\text{tight}}=r\,g_{\text{loose}}$ and taking $r\to0^+$ gives
$A_\alpha\to1$ with no CPT content at all.

Consequently, the observed $|A_\alpha|\approx1.000$ across $g_Ag_A$ pairs in
SPINDEP is \emph{consistent with}, but not \emph{evidence for}, a non-zero
$b_\mu$: the same numerical pattern is exactly what independent bounds of
very different sensitivity would produce even under exact CPT symmetry.
Distinguishing the two would require comparing bounds of comparable
sensitivity, or using actual signed measured values with symmetric
uncertainties, rather than one-sided upper limits.

\subsection{Matching to Dobrescu--Mocioiu Potentials}

The DM catalogue describes two-body potentials from boson exchange between
particles 1 and 2. The single-particle Hamiltonian of
Eq.~\eqref{eq:HNR-bi-matter} is the same Zeeman-like structure that an
axial-vector-boson exchange produces at \emph{each} vertex (cf.\ the
axial-vector/axial-vector coupling in Cong et al.\ 2025, Eq.~47). Promoting
to the two-body potential -- matching the exchange of the same axial-vector
mediator between \emph{both} fermions -- gives the spin-spin form
\begin{equation}
  V_2(r) = -g_A^{(1)}g_A^{(2)}\,\frac{1}{4\pi}\,
  (\boldsymbol\sigma_1\cdot\boldsymbol\sigma_2)\,\frac{e^{-r/\lambda}}{r},
  \qquad \lambda = 1/m_\phi. \label{eq:V2-final}
\end{equation}
The coupling constant $g_A$ is identified with the $b_\mu$ coefficient
through the matching condition $g_A\Leftrightarrow|b_\mu|/m_f$ for the
relevant fermion mass $m_f$. The spatial $b_i$ term maps to $V_2$
(spin--spin) at leading order: each fermion contributes one
$-\mathbf b\cdot\boldsymbol\sigma$ vertex, and the product of the two
single-particle Hamiltonians yields the $\boldsymbol\sigma_1\cdot\boldsymbol\sigma_2$
structure.

\begin{table}[h]
\centering
\begin{tabular}{@{}lllll@{}}
\toprule
$b_\mu$ component & NR Hamiltonian & DM potential & Order in $1/m$ & Notes \\
\midrule
$b_i$ (spatial)  & $-\mathbf b\cdot\boldsymbol\sigma$          & $V_2$      & $m^0$ (leading) & Spin--spin ($\boldsymbol\sigma_1\cdot\boldsymbol\sigma_2$); both spins couple \\
$b_0$ (temporal) & $+b_0(\boldsymbol\sigma\cdot\mathbf p)/m$   & $V_7, V_8$ & $m^{-1}$        & Velocity-dependent spin coupling \\
\bottomrule
\end{tabular}
\caption{Summary of $b_\mu$ matchings to Dobrescu--Mocioiu potentials.}
\end{table}

\subsection{Physical Consequences and SPINDEP Implications}

\textbf{All $g_Ag_A$ pairs show $|A_\alpha|\approx1$.} This is
\emph{consistent with} a $b_\mu$ sign flip, but -- per
Section~\ref{sec:bmu-asymmetry} -- is equally well explained by a
sensitivity gap between the matter- and antimatter-sector bounds, with no
CPT violation required. The two explanations cannot be distinguished from
the asymmetry value alone; doing so would require checking whether the
matter and antimatter bounds being compared are of comparable experimental
precision. Small deviations from $1$ in some pairs reflect curve-curvature
effects in the per-point uncertainty model, not necessarily physics.

\textbf{The $g_sg_s$ pair at $|A_\alpha|=0.873$.} The $g_sg_s$ potential does
not couple to $b_\mu$ at leading order (scalar coupling requires $c_{\mu\nu}$,
which is CPT-even), so a value below $1$ is consistent with no CPT-odd
contribution. This is the same type of argument as the $g_Ag_A$ case above,
however: a lower or higher $|A_\alpha|$ is equally compatible with a smaller
or larger sensitivity gap between the two experiments being compared,
independent of the CPT status of the dominant operator.

\textbf{Observable consequences.} A measurement of $b_i\neq0$ for any
fermion would appear as a direction-dependent shift in atomic energy levels
(sidereal variation). SPINDEP constraints bound the effective $b_i$ from the
ratio of matter to antimatter coupling bounds.

\section{Foldy--Wouthuysen Derivation of the \texorpdfstring{$H_{\mu\nu}$}{H\_mu\_nu} Term}
\label{sec:Hmunu}

\subsection{Overview}

The SME coefficient $H_{\mu\nu}$ is a rank-2 antisymmetric tensor coupling,
\begin{equation}
  \mathcal L_H = -\frac12 H_{\mu\nu}\,\bar\psi\,\sigma^{\mu\nu}\,\psi,
  \qquad \sigma^{\mu\nu}=\frac{i}{2}[\gamma^\mu,\gamma^\nu].
\end{equation}
Unlike $b_\mu$, $H_{\mu\nu}$ is CPT-even but Lorentz-odd: it does not change
sign under CPT, so matter and antimatter couple identically at tree level.
Its primary experimental signature is an anisotropy (direction-dependence)
rather than a matter--antimatter asymmetry. $H_{\mu\nu}$ has two distinct
sectors -- the magnetic-like spatial components $H_{ij}$ and the
electric-like mixed components $H_{0i}$ -- which generate different DM
potentials at different orders in the $1/m$ expansion.

\paragraph{Key result.} $H_{ij}$ generates $H_{NR}=-\mathcal H_B\cdot\boldsymbol\sigma$
where $\mathcal H_B^k=\tfrac12\varepsilon^{ijk}H_{ij}$, matching to $V_3$
(dipole--dipole). $H_{0i}$ generates
$H_{NR}=-\tfrac1m\boldsymbol\sigma\cdot(\mathbf p\times\mathbf H_E)$, matching
to $V_7$ (spin--velocity). Since $H_{\mu\nu}$ is CPT-even, it predicts no
$A_\alpha\neq0$ signature from the sign-flip mechanism (same sign for matter
and antimatter); its experimental signature is sidereal variation, not
matter--antimatter asymmetry.

\subsection{The SME Lagrangian Term}

The tensor $H_{\mu\nu}=-H_{\nu\mu}$ is antisymmetric and real, with 6
independent components: 3 magnetic-like, $H_{ij}=\tfrac12\varepsilon_{ijk}B_H^k$,
and 3 electric-like, $H_{0i}=E_H^i$.

Under CPT, $\bar\psi\sigma^{\mu\nu}\psi\to+\bar\psi\sigma^{\mu\nu}\psi$ (a
CPT-even bilinear), so $\mathcal L_H\to-\tfrac12H_{\mu\nu}\bar\psi\sigma^{\mu\nu}\psi$,
the same sign as the original: $H_{\mu\nu}$ is CPT-even. As a fixed
background tensor it explicitly breaks Lorentz invariance but not CPT. Under
$C$ alone, $\bar\psi^c\sigma^{\mu\nu}\psi^c=+\bar\psi\sigma^{\mu\nu}\psi$
(a $C$-even tensor current), so the antiparticle coupling is the
\emph{same} sign as the particle: $H_{NR}^{\bar f}=H_{NR}^f$. This means
$H_{\mu\nu}$ predicts $A_\alpha=0$ in the absence of sensitivity differences.

\subsection{FW Transformation of \texorpdfstring{$H_{\mu\nu}$}{H\_mu\_nu}}

\paragraph{Magnetic-like components $H_{ij}$.} The relevant combination is
$\gamma^0\sigma^{ij}$ where $\sigma^{ij}=\tfrac{i}{2}[\gamma^i,\gamma^j]$. In
the Dirac representation,
\begin{equation}
  \gamma^0\sigma^{ij} = \varepsilon^{ijk}\Sigma^k
  = \mathrm{diag}\!\left(\varepsilon^{ijk}\sigma^k,\varepsilon^{ijk}\sigma^k\right),
\end{equation}
an even (block-diagonal) operator. In the upper components, contracting with
$H_{ij}$,
\begin{equation}
  H_{H_{ij}}^\uparrow = -\frac12 H_{ij}\varepsilon^{ijk}\sigma^k
  = -\mathcal H_B^k\sigma^k
  \;\Longrightarrow\;
  \boxed{H_{NR}(H_{ij}) = -\boldsymbol{\mathcal H}_B\cdot\boldsymbol\sigma.}
\end{equation}
This is the standard spin--magnetic-field coupling, with the SME background
$\boldsymbol{\mathcal H}_B$ playing the role of a magnetic field. It
generates the dipole--dipole potential $V_3$ via the tensor exchange
propagator.

\paragraph{Electric-like components $H_{0i}$.} The mixed components $H_{0i}$
contribute through $\gamma^0\sigma^{0i}=\gamma^0\times\tfrac{i}{2}[\gamma^0,\gamma^i]$.
In the Dirac representation,
\begin{equation}
  \gamma^0\sigma^{0i} = \frac12\gamma^0[\gamma^0,\gamma^i] = -\alpha^i
  \quad (\text{off-diagonal}),
\end{equation}
an odd operator mixing upper and lower spinor components. At order $m^0$ it
vanishes; the FW procedure generates a contribution at order $m^{-1}$,
\begin{equation}
  H_{NR}(H_{0i}) = -\frac1m\,\boldsymbol\sigma\cdot(\mathbf p\times\mathbf H_E),
  \qquad H_E^i \equiv H_{0i},
\end{equation}
a spin--orbit-type coupling that generates DM potential $V_7$.

\subsection{Two-Body Potential Matching}

The magnetic-like $H_{ij}$ term generates the dipole--dipole potential
\begin{equation}
  V_3 = \frac{g^2}{4\pi}\left[\boldsymbol\sigma_1\cdot\boldsymbol\sigma_2
  - 3(\boldsymbol\sigma_1\cdot\hat{\mathbf r})(\boldsymbol\sigma_2\cdot\hat{\mathbf r})\right]
  \frac{e^{-r/\lambda}}{r},
\end{equation}
where $g$ is identified with the effective $H_{ij}$ coupling; both particles
must carry a spin coupling for $V_3$ to contribute. The electric-like
$H_{0i}$ term generates the spin--velocity potential
\begin{equation}
  V_7 = \frac{g^2}{4\pi}\,\boldsymbol\sigma_1\cdot(\mathbf v\times\hat{\mathbf r})\,
  \frac{e^{-r/\lambda}}{r},
\end{equation}
which is velocity-dependent and requires experimental sensitivity to
relative motion between the test masses.

\begin{table}[h]
\centering
\begin{tabular}{@{}lllll@{}}
\toprule
$H_{\mu\nu}$ sector & NR Hamiltonian & DM potential & Order in $1/m$ & Signature \\
\midrule
$H_{ij}$ (magnetic) & $-\boldsymbol{\mathcal H}_B\cdot\boldsymbol\sigma$ & $V_3$ & $m^0$ (leading) & Sidereal variation in spin precession \\
$H_{0i}$ (electric) & $-\tfrac1m\boldsymbol\sigma\cdot(\mathbf p\times\mathbf H_E)$ & $V_7$ & $m^{-1}$ & Velocity/direction-dependent force \\
\bottomrule
\end{tabular}
\caption{Summary of $H_{\mu\nu}$ matchings to Dobrescu--Mocioiu potentials.}
\end{table}

\section{Foldy--Wouthuysen Derivation of the \texorpdfstring{$d_{\mu\nu}$}{d\_mu\_nu} Term}
\label{sec:dmunu}

\subsection{Overview}

The SME coefficient $d_{\mu\nu}$ is a CPT-even, Lorentz-odd rank-2 tensor
coupling involving $\gamma_5$: $\mathcal L_d=d_{\mu\nu}\bar\psi\gamma_5\sigma^{\mu\nu}\psi$.
It is the axial-tensor partner of $H_{\mu\nu}$. Despite being CPT-even (same
sign for matter and antimatter), it generates a rich phenomenology because
different components produce NR Hamiltonians at different orders in $1/m$,
spanning potentials $V_2$, $V_7$, and $V_8$. The term is particularly
relevant for electron--nucleon scattering experiments: $d_{i0}$ produces a
large (order $m^1$) but momentum-independent coupling to $\sigma^i$, $d_{ij}$
produces a velocity-dependent coupling, and $d_{00}=0$ identically by
antisymmetry.

\begin{quote}
\textbf{Important caveat -- unresolved as of this writing.} While completing
the executable symbolic-derivation notebooks for $b_\mu$ and $H_{\mu\nu}$
(Sections~\ref{sec:bmu} and~\ref{sec:Hmunu}), a genuine Dirac-algebra sign
error was found in the $b_\mu$ derivation and fixed only after actually
running the symbolic computation, rather than trusting the hand-derived
prose. Applying that same numeric check to the electric-axial $d_{i0}$ term
below gives an \emph{even}, $O(m^0)$, \emph{imaginary}-coefficient operator
-- not the real, $O(m^1)$, mass-enhanced $+d_{i0}\,m\,\sigma^i$ claimed
below. The claimed order-$m^1$ enhancement does not fall out of the same
direct $\gamma^0\Gamma$ projection that correctly reproduced the $b_\mu$ and
$H_{\mu\nu}$ results, and re-deriving it properly needs more careful
multi-step work than has been completed so far. \textbf{The key results
below should therefore be treated as unverified} pending a dedicated,
executed symbolic derivation analogous to the ones already completed for
$b_\mu$ and $H_{\mu\nu}$; the $d_{i0}\to V_2$ mass-enhancement claim in
particular should not be relied upon until that verification is done.
\end{quote}

\paragraph{Key results (unverified -- see caveat above).}
\begin{align}
  d_{i0}&: \; H_{NR} = +d_{i0}\,m\,\sigma^i \quad (\text{order } m^1,\text{ large but static}) \;\to V_2 \\
  d_{ij}&: \; H_{NR} = +d_{ij}\,p^j\,\sigma^i \quad (\text{order } m^0 \text{ in velocity}) \;\to V_7, V_8 \\
  d_{00}&: \; H_{NR} = +d_{00}\,\boldsymbol\sigma\cdot\mathbf p \quad (\text{order } m^0 \text{ in momentum}) \;\to V_8
\end{align}
$d_{\mu\nu}$ is CPT-even: no sign flip for antimatter is predicted, and
$A_\alpha$ is expected near $0$ subject to the same sensitivity-gap caveat
developed in Section~\ref{sec:bmu-asymmetry}.

\subsection{The SME Lagrangian Term}

The $d_{\mu\nu}$ Lagrangian is
\begin{equation}
  \mathcal L_d = d_{\mu\nu}\,\bar\psi\,\gamma_5\,\sigma^{\mu\nu}\,\psi,
  \qquad \sigma^{\mu\nu}=\frac{i}{2}[\gamma^\mu,\gamma^\nu].
\end{equation}
$d_{\mu\nu}=-d_{\nu\mu}$ is real and antisymmetric with 6 independent
components: three $d_{i0}$ (electric-axial) and three $d_{ij}$
(magnetic-axial); $d_{00}=0$ identically by antisymmetry. (In some SME
references $d_{\mu\nu}$ absorbs additional factors; this chapter follows the
convention of Kostelecký and Lane, 1999.)

Under CPT, $\bar\psi\gamma_5\sigma^{\mu\nu}\psi\to+\bar\psi\gamma_5\sigma^{\mu\nu}\psi$
(the combination $\gamma_5\sigma^{\mu\nu}$ is CPT-even), so $\mathcal L_d$ is
CPT-even. Under $C$, $\bar\psi^c\gamma_5\sigma^{\mu\nu}\psi^c=+\bar\psi\gamma_5\sigma^{\mu\nu}\psi$
(same sign): the $d_{\mu\nu}$ contribution to the effective coupling is
therefore the same for matter and antimatter.

\subsection{FW Derivation by Component}

\paragraph{Electric-axial components $d_{i0}$.} The relevant combination is
$\gamma^0\gamma_5\sigma^{i0}=\gamma^0\gamma_5\,\tfrac{i}{2}[\gamma^i,\gamma^0]$.
In the Dirac representation, as claimed (subject to the caveat above),
\begin{equation}
  \gamma^0\gamma_5\sigma^{i0} = +m\Sigma^i = \mathrm{diag}(+m\sigma^i,+m\sigma^i),
\end{equation}
an even operator at order $m^1$, giving
$H_{NR}(d_{i0})=+d_{i0}\,m\,\sigma^i$. This term is large (enhanced by $m$)
but does not produce a spatial gradient and so does not contribute to a
force law directly; in a two-body interaction it would contribute to $V_2$
as a background spin-polarising field, with correspondingly tight
experimental bounds from torsion-balance experiments.

\paragraph{Magnetic-axial components $d_{ij}$.} For the spatial
antisymmetric components, the relevant combination is
$\gamma^0\gamma_5\sigma^{ij}=\tfrac{i}{2}\gamma^0\gamma_5[\gamma^i,\gamma^j]$.
Expanding in the Dirac representation, one obtains an odd matrix at zeroth
order that becomes even at order $m^{-1}$ after the FW transformation,
involving the momentum operator:
\begin{equation}
  H_{NR}(d_{ij}) = +d_{ij}\,p^j\,\sigma^i \quad (\text{summed over } i,j),
\end{equation}
a velocity-dependent coupling to the spin, generating DM potentials $V_7$
and $V_8$.

\paragraph{Pure temporal $d_{00}=0$.} By antisymmetry of $d_{\mu\nu}$,
$d_{00}=0$ identically, so there is no NR contribution from this component.

\subsection{Two-Body Potentials and Matching}

The large $d_{i0}m\sigma^i$ coupling would contribute to $V_2$ in the
two-body potential as the spin-polarised vertex. Since $H_{NR}(d_{i0})$ has
the same $\sigma^i$ (Zeeman-like) structure as the $b_i$ case
(Section~\ref{sec:bmu-application}), the two-body potential it generates
would likewise be spin--spin, not monopole--dipole:
\begin{equation}
  V_2\text{-contribution} \propto (d_{i0}m)^{(1)}(d_{i0}m)^{(2)}\times
  (\boldsymbol\sigma_1\cdot\boldsymbol\sigma_2)\,\frac{e^{-r/\lambda}}{r},
\end{equation}
again subject to the verification caveat above. The large $m$ factor means
even very small $d_{i0}$ coefficients would produce potentially observable
effects; current bounds from atomic magnetometry constrain
$|d_{i0}|<10^{-25}$~GeV (Heckel et al., 2008). The velocity-dependent
$d_{ij}$ term would generate
\begin{align}
  V_7 &\propto d_{ij}\,\sigma^i\,v^j\,\frac{e^{-r/\lambda}}{r}, \\
  V_8 &\propto d_{ij}\left\{(\boldsymbol\sigma\cdot\mathbf v),\,\frac{e^{-r/\lambda}}{r}\right\}
  \quad (\text{anticommutator}).
\end{align}

\begin{table}[h]
\centering
\begin{tabular}{@{}lllll@{}}
\toprule
$d_{\mu\nu}$ component & NR Hamiltonian & DM potential(s) & Order in $1/m$ & Notes \\
\midrule
$d_{i0}$ (electric-axial) & $+d_{i0}\,m\,\sigma^i$ & $V_2$ & $m^1$ (large) & Enhanced by fermion mass; unverified \\
$d_{ij}$ (magnetic-axial) & $+d_{ij}\,p^j\,\sigma^i$ & $V_7, V_8$ & $m^0$ in velocity & Velocity-dependent; needs moving sources \\
$d_{00}$ & $=0$ by antisymmetry & None & --- & Identically zero \\
\bottomrule
\end{tabular}
\caption{Summary of $d_{\mu\nu}$ matchings to Dobrescu--Mocioiu potentials (unverified, see caveat above).}
\end{table}

\section{SME Coefficient--Dobrescu--Mocioiu Potential Matching Table (Master Reference)}
\label{sec:master}

This section synthesises the three FW derivations above into a single
lookup table, matching each SME coefficient onto the corresponding DM
potential(s) from the complete $V_1$--$V_{16}$ catalogue given in
Chapter~2. Its primary use is to interpret measured values of $A_\alpha$ in
terms of bounds on specific SME coefficients, and to predict which DM
potential a given SME background field generates in a two-body interaction.
Notation follows Chapter~2 throughout: $\boldsymbol\sigma_{1,2}$ for the
Pauli matrices of particles 1 and 2, $\hat{\mathbf r}$ for the unit
separation vector, $\lambda=1/m_\phi$ for the interaction range, and
$\mathcal H_B^k=\tfrac12\varepsilon^{ijk}H_{ij}$ for the dual magnetic-like
vector built from $H_{\mu\nu}$.

\subsection{Main SME\texorpdfstring{$\to$}{->}DM Potential Matching Table}

Table~\ref{tab:master} gives the complete mapping from SME coefficient to NR
Hamiltonian to DM potential(s), with CPT/Lorentz properties and the naive
predicted $A_\alpha$ signature. A caveat on the prediction column is
essential: it is tempting to substitute an exact, signed CPT-odd relation
($g_{\bar f}=-g_f$) directly into the SPINDEP formula
$A_\alpha=(g_f-g_{\bar f})/(g_f+g_{\bar f})$ and conclude $|A_\alpha|\to1$.
That substitution does not hold -- the denominator vanishes and the
expression diverges instead, as derived in Section~\ref{sec:bmu-asymmetry}.
More importantly, SPINDEP's actual inputs are independent, always-positive
experimental upper bounds, not a signed measured coupling, so the signed
substitution does not describe what is actually computed. Comparing two
positive bounds of very different tightness drives $|A_\alpha|\to1$
regardless of CPT status -- a sensitivity-gap effect. The predictions below
should therefore be read as \emph{consistent with}, not \emph{caused by},
the stated CPT parity.

\begin{table}[h]
\centering
\small
\begin{tabular}{@{}lllllll@{}}
\toprule
SME coeff. & CPT & Lorentz & NR Hamiltonian & DM potential(s) & $1/m$ order & SPINDEP $A_\alpha$ prediction \\
\midrule
$b_i$ & Odd & Odd & $-\mathbf b\cdot\boldsymbol\sigma$ (matter), $+\mathbf b\cdot\boldsymbol\sigma$ (antimatter) & $V_2$ & $m^0$ & Consistent with $|A_\alpha|$ near 1, but equally consistent with a sensitivity gap \\
$b_0$ & Odd & Odd & $+b_0(\boldsymbol\sigma\cdot\mathbf p)/m$ & $V_7, V_8$ & $m^{-1}$ & Same caveat as $b_i$ \\
$H_{ij}$ & Even & Odd & $-\boldsymbol{\mathcal H}_B\cdot\boldsymbol\sigma$ & $V_3$ & $m^0$ & $A_\alpha\approx0$ if bounds comparably sensitive; sensitivity gap can still give $|A_\alpha|$ near 1 \\
$H_{0i}$ & Even & Odd & $-\tfrac1m\boldsymbol\sigma\cdot(\mathbf p\times\mathbf H_E)$ & $V_7$ & $m^{-1}$ & Same caveat as $H_{ij}$ \\
$d_{i0}$ & Even & Odd & $+d_{i0}\,m\,\sigma^i$ & $V_2$ & $m^1$ (!) & Same caveat as $H_{ij}$; unverified (Sec.~\ref{sec:dmunu}) \\
$d_{ij}$ & Even & Odd & $+d_{ij}\,p^j\,\sigma^i$ & $V_7, V_8$ & $m^0$ (vel.) & Same caveat as $H_{ij}$ \\
$d_{00}$ & Even & Even & $=0$ by antisymmetry & None & --- & No contribution \\
\bottomrule
\end{tabular}
\caption{Master SME$\to$DM matching table.}
\label{tab:master}
\end{table}

Because every SPINDEP $A_\alpha$ value is a ratio of two independent
one-sided bounds, an observed value near $\pm1$ can arise either from a
genuine CPT-odd signal \emph{or} from nothing more than one experiment being
much more sensitive than the other. Distinguishing the two requires
independent information about the relative precision of the matter- and
antimatter-sector measurements -- it cannot be read off $A_\alpha$ alone.
This applies uniformly to every row of Table~\ref{tab:master}, CPT-odd or
CPT-even.

\subsection{Connecting to SPINDEP Results}

\subsubsection{Validated \texorpdfstring{$g_Ag_A$}{gAgA} Pairs}

\begin{table}[h]
\centering
\small
\begin{tabular}{@{}lllll@{}}
\toprule
SPINDEP pair & Coupling & DM potential & Dominant SME coeff. & Observed $|A_\alpha|$ -- consistent with, not proof of \\
\midrule
$g_sg_s\cdot$UNKNOWN$\cdot ee$ & Scalar-scalar & $V_1$ (UNKNOWN) & $c_{\mu\nu}$ (CPT-even) & $0.873$: consistent with a sensitivity-gap-dominated CPT-even channel \\
$g_Ag_A\cdot V_1\cdot ep$ & Axial-axial & $V_2$ (spin-spin) & $b_\mu$ (CPT-odd) & $0.9998$: consistent with a CPT-odd sign flip, equally consistent with a pure sensitivity gap \\
$g_Ag_A\cdot V_2\cdot ee$ ($\times5$) & Axial-axial & $V_2$ (spin-spin) & $b_\mu$ (CPT-odd) & $0.954$--$1.000$: same caveat \\
\bottomrule
\end{tabular}
\end{table}

\subsubsection{Interpretation of the \texorpdfstring{$g_sg_s$}{gsgs} Result}

The $g_sg_s\cdot$UNKNOWN$\cdot ee$ pair shows $|A_\alpha|=0.873$, lower than
all $g_Ag_A$ pairs. Two points are worth noting. First, the scalar--scalar
($g_sg_s$) coupling does not appear in the minimal SME at dimension 4; its
primary contribution comes from $c_{\mu\nu}$, a CPT-even coefficient. A
value below the $g_Ag_A$ pairs is \emph{consistent with} a CPT-even channel
and a sizeable sensitivity gap (the matter constraint used here is several
orders of magnitude tighter than the antimatter-sector positronium
constraint it is compared against) -- but the $g_Ag_A$ pairs' higher values
are \emph{equally} explainable by a sensitivity gap alone, so the comparison
between rows in Table~\ref{tab:master} cannot by itself distinguish
CPT-odd from CPT-even-with-a-larger-sensitivity-gap. Second, the potential
label UNKNOWN indicates the parser could not match the filename to a $V_n$
token; the likely candidate is $V_1$ (monopole--monopole, scalar exchange)
given the $g_sg_s$ coupling. $V_1$ has no spin structure and would produce
$A_\alpha=0$ for a CPT-symmetric world \emph{if} the compared bounds were of
comparable sensitivity.

\subsubsection{CPT Rule for All SME Coefficients}

The general rule, derived from the FW analysis above, is
\begin{align}
  H_{NR}^{\bar f}\!\left(X^{\text{CPT-odd}}\right)  &= -H_{NR}^{f}\!\left(X^{\text{CPT-odd}}\right), \\
  H_{NR}^{\bar f}\!\left(X^{\text{CPT-even}}\right) &= +H_{NR}^{f}\!\left(X^{\text{CPT-even}}\right).
\end{align}
This is a statement about a single, hypothetical \emph{signed} coupling and
its antiparticle counterpart. It does \emph{not}, by itself, imply anything
about the ratio of two independent experimental upper bounds. For a
CPT-odd coefficient such as $b_\mu$, an exact signed relation
$g_{\bar a}=-g_m$ makes the SPINDEP formula
$A_\alpha=(g_m-g_{\bar a})/(g_m+g_{\bar a})$ undefined or divergent, not a
bounded $|A_\alpha|=1$; what actually produces $|A_\alpha|\to1$ in real
SPINDEP output is two independent positive bounds of very different
tightness, which happens regardless of CPT parity. For a CPT-even
coefficient ($H_{\mu\nu}$, $d_{\mu\nu}$, $c_{\mu\nu}$), $g_{\bar a}=+g_m$
gives $A_\alpha=0$ only if $g_m$ and $g_{\bar a}$ are numerically equal --
again a statement about a hypothetical signed measurement, not about
independent one-sided bounds, which will show whatever sensitivity gap
exists between the two experiments regardless of the true CPT-even
relation.

In short: $A_\alpha$ computed from upper bounds cannot currently distinguish
CPT violation from an experimental sensitivity gap. Doing so would require
either bounds of demonstrably comparable sensitivity, or actual
central-value measurements with symmetric errors instead of one-sided
limits -- a limitation this thesis returns to when interpreting the full
constraint database in Chapter~4.

\subsection{Additional SME Coefficients (for completeness)}

\paragraph{$a_\mu$ (CPT-odd, Lorentz-odd).} The term
$\mathcal L_a=a_\mu\bar\psi\gamma^\mu\psi$ produces $H_{NR}=-a_0$ at order
$m^0$ (a scalar energy shift) and $-a_iv^i/m$ at order $m^{-1}$
(velocity-dependent). It does not generate a spin-dependent potential at
leading order and therefore does not appear in the DM catalogue directly. It
is CPT-odd, so antimatter picks up the opposite sign -- relevant for clock
comparisons (ALPHA, BASE) but not for torsion-balance spin-dependent
constraints.

\paragraph{$c_{\mu\nu}$ (CPT-even, Lorentz-odd).} The term
$\mathcal L_c=c_{\mu\nu}\bar\psi\gamma^\mu i\partial^\nu\psi$ generates
momentum-dependent (direction-dependent) energy shifts. It is CPT-even, and
in the NR limit contributes a direction-dependent kinetic-energy
modification that mimics a preferred-frame effect. For spin-dependent
interactions it generates $V_2$-type couplings at subleading order; the
$g_sg_s$ sector may be sensitive to $c_{\mu\nu}$ if the mediator couples
through the scalar current.

\paragraph{$e_\mu, f_\mu$ (CPT-odd, higher dimension).} Dimension-4+
operators $e_\mu$ and $f_\mu$ are CPT-odd but velocity-dependent (odd under
$C$ but even under $P$). They contribute to $V_8$ and higher DM potentials.
They are not yet constrained by SPINDEP pairs but should become accessible
as the $g_Vg_V$ and $g_pg_p$ pairs are analysed with an updated parser.

\section{Chapter Summary}

This chapter derived the non-relativistic Hamiltonians generated by $b_\mu$,
$H_{\mu\nu}$, and $d_{\mu\nu}$ via the Foldy--Wouthuysen transformation and
matched each to the Dobrescu--Mocioiu potential catalogue. The $b_\mu$ and
$H_{\mu\nu}$ derivations were verified by executable symbolic computation;
the $d_{\mu\nu}$ derivation is flagged as unverified pending further work
and should not yet be cited as an established result. Throughout, a single
interpretive caveat applies uniformly to every coefficient: because SPINDEP
compares independent, one-sided experimental upper bounds rather than
signed measured couplings, an observed asymmetry $|A_\alpha|$ near $1$ is
consistent with, but not evidence for, CPT violation -- the same pattern
arises from a sensitivity gap between the matter- and antimatter-sector
experiments regardless of the true CPT status of the underlying physics.
Chapter~4 applies this framework to the full compiled constraint database.

\ifSubfilesClassLoaded{
\begin{thebibliography}{99}

\bibitem{Colladay1997}
D.~Colladay and V.~A. Kostelek\'{y},
\textit{CPT violation and the standard model},
Phys.\ Rev.\ D \textbf{55}, 6760 (1997).

\bibitem{Colladay1998}
D.~Colladay and V.~A. Kostelek\'{y},
\textit{Lorentz-violating extension of the standard model},
Phys.\ Rev.\ D \textbf{58}, 116002 (1998).

\bibitem{Dobrescu2006}
B.~A. Dobrescu and I.~Mocioiu,
\textit{Spin-dependent macroscopic forces from new particle exchange},
J.\ High Energy Phys.\ \textbf{11}, 005 (2006).

\bibitem{KosteleckySamuel1989}
V.~A. Kostelek\'{y} and S.~Samuel,
\textit{Spontaneous breaking of Lorentz symmetry in string theory},
Phys.\ Rev.\ D \textbf{39}, 683 (1989).

\bibitem{KosteleckyLane1999}
V.~A. Kostelek\'{y} and C.~D. Lane,
\textit{Nonrelativistic quantum Hamiltonian for Lorentz violation},
Phys.\ Rev.\ D \textbf{60}, 116010 (1999).

\bibitem{KosteleckyMewes2001}
V.~A. Kostelek\'{y} and M.~Mewes,
\textit{CPT violation and the standard model},
Phys.\ Rev.\ D \textbf{66}, 056005 (2001).

\bibitem{BaileyKostelecky2006}
Q.~G. Bailey and V.~A. Kostelek\'{y},
\textit{Signals for Lorentz violation in post-Newtonian gravity},
Phys.\ Rev.\ D \textbf{74}, 045001 (2006).

\bibitem{Foldy1950}
L.~L. Foldy and S.~A. Wouthuysen,
\textit{On the Dirac theory of spin-$\frac{1}{2}$ particles and its
non-relativistic limit},
Phys.\ Rev.\ \textbf{78}, 29 (1950).

\bibitem{Heckel2008}
B.~R. Heckel, E.~G. Adelberger, C.~E. Cramer, T.~S. Cook, S.~Schlamminger,
and U.~Schmidt,
\textit{Preferred-frame and CP-violation tests with polarized electrons},
Phys.\ Rev.\ D \textbf{78}, 092006 (2008).

\bibitem{Fadeev2022}
P.~Fadeev, F.~Ficek, M.~G. Kozlov, D.~Budker, and V.~V. Flambaum,
\textit{Constraints on exotic spin-dependent interactions between electrons from
helium fine-structure spectroscopy},
Phys.\ Rev.\ A \textbf{105}, 022812 (2022).

\bibitem{Cong2025}
L.~Cong \textit{et al.},
\textit{Spin-dependent exotic interactions},
Rev.\ Mod.\ Phys.\ \textbf{97}, 025005 (2025).

\end{thebibliography}
}{}

\end{document}
